% Translator-owned audit and supplemental proofs for the initial part of section 13.
\head{第13节初段：定义、维数归约与初等邻域构造}

\noindent\textbf{实际范围与证据。}
本模块逐页查看法文手稿\texttt{EGA5\_francais.pdf}物理第74--77页、英文编辑稿\texttt{EGAVaiello.pdf}物理及印刷第49--51页的页面影像，并放大核对法文第76页中心与参数空间一段。法文第78页另作下一接点检查，未计入本模块译文。有效范围从第74页第13节标题开始，止于第77页末“C'est le voisinage cherch\'e de $x$ contenu dans $X-Y$”。英文对应止于第51页第一段末。其后的引理13.3证明交由下一模块，不重译、不重复输入。

本次接续基线为公开版本\texttt{99cd048}中的实际文件：
\par\noindent
\path{current/EGA_v20/continuations/v20_section13_inherited.tex}
\par\smallskip
本模块对应其中旧D17-M046至D17-M060。原文件保持不变；这15条旧注的原文另逐条保存于本报告的历史附录。旧注将“底本”笼统指向英文稿的地方，不能直接改当对法文稿的判断。本文以下标作“补证”的各节均属于本轮译者的数学核验，不是对作者删去或没有写出的证明的复原。

\subsection*{一、逐段对应及实际修订}

\begin{longtable}{@{}p{.12\linewidth}p{.17\linewidth}p{.17\linewidth}p{.42\linewidth}@{}}
\toprule
源码标记 & 法文物理页 & 英文物理页 & 实际内容与接点\\
\midrule
\endhead
INI-P01 & 74上 & 49上 & 定义13.1、复合定义、英文编辑脚注63。\\
INI-P02 & 74中 & 49中 & 定理13.2、给定邻域、维数归纳与空复合提醒。\\
INI-P03 & 74末--75首 & 49下 & 曲线情形、正规射影模型、完美性及英文脚注64。\\
INI-P04 & 75上 & 49末 & 仿射、射影闭包、正规化及未填引文、英文脚注65。\\
INI-P05 & 75中下 & 50上 & 奇异闭集、超曲面边界、约化结构及扩大$Z$。\\
INI-P06 & 75末--76中 & 50中 & 固定中心的想法、引理13.3全部陈述与无限域括注。\\
INI-P07 & 76中 & 50中下 & 从引理构造邻域、横截中心$C$及非空平展$T$。\\
INI-P08 & 76末--77首 & 50末--51首 & $Q$、$\xi_0$、三个关联族及三个结构投影。\\
INI-P09 & 77中 & 51上 & 几何连通性、基底开限制、边界与初等态射。\\
INI-P10 & 77末 & 51首段末 & 中心外同构、删边界后的纤维、包含$x$的结论。\\
\bottomrule
\end{longtable}

\begin{longtable}{@{}p{.17\linewidth}p{.35\linewidth}p{.39\linewidth}@{}}
\toprule
本版注号 & 源页证据 & 错误归属与本版处理\\
\midrule
\endhead
INI001 & 法74定义写$X'\to X$；英49写$X'\to Y$。 & 法稿原错；正文用$Y$，承认英文此前已正确。\\
INI002、005 & 法74有旧术语删划及空复合编写提醒。 & 保留有效定义与作者提醒；不复活被划去词。\\
INI003、007 & 法74定理未限有理点；法75正规化留“r\'ef.”。 & 不擅加全局有限型假设；从局部光滑性补足有限型归约，历史引文仍不冒充已定位。\\
INI004 & 两稿都误用“初等”作任意维数的最终邻域。 & 订为“多重初等”；法稿句末字未辨，保留英文疑问。\\
INI008 & 法75扭曲次数也用$n$。 & 次数改记$d$，维数$n$不动；另给边界存在的补证。\\
INI009 & 法76明确$X\supset Y\supset Z$、$\phi_n i$、quelconque。 & 包含链错误、上标误排、“some?”来自英文；法稿重复维数句照留。\\
INI010 & 法76引理陈述未写$x\in L_0$，法78证明却从经过$x$的截面开始。 & 调用处补明经过$x$的加强形式；一般闭点衔接由13.4(c)补证处理。\\
INI011 & 法76明确$C\subset L_0$；删改后有效选取句只剩避开有限集。 & 英文漏0；正文另补横截条件并给译者证明，划去旧句只作历史证据。\\
INI012 & 法76为$\xi_0$，法77为$T_Q^{n-1}$；定义$Q$时一字覆打。 & 英文漏下标或对象；覆打字不强行辨认，按固定中心及英文写$C$。\\
INI013、014 & 法77称$p$光滑，并以“Bertini”说明连通。 & 补核相对Jacobian条件；连通性引用的完整射影不可约假设已核。\\
INI015 & 法77为“$\xi_0$在$X$中的开邻域$V$”。 & 环境$X$确为法稿原错，应为$Q$；英文另错$U$及点名，分别归责。\\
INI016 & 法77只显明边界平展，直接称初等。 & 补核有限性、恒定次数、满射及射影性。\\
INI017、018 & 法77误写$Z_Q$，又说包含整个$L_0\cap X$。 & 两处均为法稿原错；改成$T_Q$及删去$Y\cup T$后的部分。\\
\bottomrule
\end{longtable}

法文第76页在中心选取处有大段删除：可辨它曾援引第7节，并对$L_0\cap X$与中心的交提出光滑、零维等要求；随后手添避开有限集及不含$x$。本版没有把删除段重排为原有效证明。正文新增的“横截”明确标为译者为使推论成立所补的条件，存在性证明另列于下。法文第74页的旧术语及第75页若干覆打旧句同样按删除层处理；无法可靠辨认的原字未据数学需要反造。

\subsection*{二、初段补证一：有限型入口、零维、正规化与边界}

\noindent\textbf{1. 光滑点提供所需有限表示邻域。}
设$O$是$x$的任意给定开邻域。$X\to\Spec(k)$在$x$处光滑，按光滑性的局部定义，存在包含$x$的仿射开集$W_0$使$W_0\to\Spec(k)$光滑且有限表示。在$W_0\cap O$内取包含$x$的主开集$W=D(g)$，则其坐标环由$W_0$的有限表示$k$-代数再局部化一个元素得到，仍有限表示。故$W$仿射、光滑、有限型，并包含在原给定邻域内。光滑性意味着几何约化；$X$几何不可约意味着这个非空开集$W$几何不可约，所以$W$几何整。

因此原定理无需额外假设整个$X$有限型或拟紧。后续射影闭包及正规化应作用于上述$W$，而不是直接对一个未经归约的任意预概形操作。可核验的定义和性质见\href{https://stacks.math.columbia.edu/tag/01V4}{Stacks Project，Tag 01V4，光滑态射的局部定义与有限表示性质}。这里将“在$x$处光滑”按态射的通常含义使用；没有把“局部环正则”擅自当成不附条件的光滑性。

\noindent\textbf{2. 零维与归纳的实际拼接。}
若$\dim W=0$，光滑有限型仿射$k$-概形$W$是有限平展的。基变换到代数闭包后，它是有限个约化点的不交和；几何不可约性迫使只有一个点。其有限平展秩为1，所以$W\simeq\Spec(k)$，结构态射用空复合表示。

若已经找到初等态射$f:U\to V$，并在$V$中找到包含$f(x)$的多重初等开邻域$V_1$，则$U_1=f^{-1}(V_1)$是$x$的开邻域。初等态射的射影光滑紧化、有限平展边界和几何连通一维纤维均对基变换保持，所以$U_1\to V_1$仍初等；与$V_1\to\Spec(k)$的初等分解复合即得所需分解。每个初等态射的几何纤维是光滑连通射影曲线删去有限个点，非空且几何不可约，因此它满射且光滑。于是若$U$几何不可约，满射像$V$也几何不可约；相对维数为1给出$\dim V=\dim U-1$。对于一般闭点，归纳必须同时携带其有限像点集合；该加强归纳在本版13.4(c)的补证中完成。

\noindent\textbf{3. 射影闭包与正规化。}
将$W$嵌入仿射空间，再取其在射影空间中的约化闭包$\overline W$。这是射影整$k$-概形，包含$W$为稠密开集。有限型$k$-概形的正规化有限，故$\nu:\widetilde W\to\overline W$有限且射影；$W$正则因而正规，所以$\nu^{-1}(W)\to W$为同构。有限正规化和射影性的明确出处是\href{https://stacks.math.columbia.edu/tag/0BXQ}{Stacks Project，Tag 0BXQ，正规化第33.27.1及33.27.2引理}。原稿未填的“r\'ef.”仍是历史引文的空缺；这里提供的是本版核验资料，不伪称它就是作者当年拟引位置。

完美域没有非平凡有限纯不可分扩张。由几何正规性的有限纯不可分扩张判据，正规$k$-概形便几何正规；参见\href{https://stacks.math.columbia.edu/tag/038N}{Stacks Project，Tag 038N}。$W$的几何整性又保证上述闭包及正规化的几何不可约性：在任意扩域后，$W$仍是稠密不可约开集；正规化在这个开集上保持同构，且有限型正规的各分量不能产生完全支持于较低维边界的新分量。也可使用正规化对可分扩域的相容性，而完美域的代数闭包扩张可分。有限型正规簇满足余维1处正则，完美域上正则点即光滑点，故非光滑点集余维至少为2。

\noindent\textbf{4. 所需纯余维1边界的存在。}
现在记正规射影几何整紧化为$X$，给定开邻域为$O$，其非光滑闭集为$Z_0$。令
\[
B=(X\setminus O)\cup Z_0.
\]
这是一个不含$x$的闭集。在固定的射影嵌入下，射影拓扑的齐次主开集构成基，因此存在正次数齐次式$h$，满足$h(x)\ne0$，且$h$在$B$上消失。必要时将次数乘高，不改变这两个条件。$h|_X$不是零截面；$X$整，故其零点为非零因子定义的有效Cartier除子。有限型整簇上的主理想定理及维数公式表明，它的每个不可约分量维数为$n-1$。令$Y$为其支集所赋的约化闭子概形，即得$B\subset Y$、$x\notin Y$、$X\setminus Y\subset O$及$Y$纯维$n-1$。

因为$k$完美，约化有限型$Y$的每个不可约分量都有非空光滑开集；其非光滑点集闭且维数至多$n-2$。令$Z$为$Z_0$与该非光滑点集的并，赋予约化闭结构，即得原段所需的$Z\subset Y$及维数界。这个论证说明原段的边界选择可行，但不把论证文字并入作者正文。

\subsection*{三、初段补证二：曲线边界及完美性确实不能无条件删去}

\noindent\textbf{1. 完美域上的曲线归约。}
先按补证一取光滑仿射几何整曲线$W$。将其嵌入射影闭包并正规化，得到正则射影曲线$C$，且$W$为$C$的开子概形。由于$k$完美，$C$光滑；$C$几何不可约，所以几何连通。补集$C\setminus W$为有限个闭点，赋约化结构后是有限平展$k$-概形。它非空：若为空，则正维$C$同时仿射和固有，而仿射固有$k$-概形有限，矛盾。因而这个边界到$\Spec(k)$满射且次数恒定，$W\to\Spec(k)$满足定义13.1。先取仿射邻域已解决旧D17-M050所指出的空边界问题，不必向原证明暗中增加一个未经选取的删点。

\noindent\textbf{2. 不完美域上的明确反例。}
以下例子是译者构造及核验，不是原文提供的例子。取
\[
k=\mathbb F_3(t),\qquad
C=V(Y^2Z-X^3-tZ^3)\subset\mathbf P^2_k.
\]
域$k$无限且不完美。记仿射图$Z=1$上的坐标环为
\[
A=k[x,y]/(y^2-x^3-t).
\]
多项式$x^3+t$在$k[x]$中不可约：否则$-t$在$k$中有三次根，而$t$在$t=0$处的赋值为1，不可能是一个元素的三次方的赋值。于是$(y)$是$A$的极大理想，剩余域为纯不可分三次扩张$k[x]/(x^3+t)$；记相应闭点为$q$。

在代数闭包中取$\alpha^3=-t$，则仿射方程变成
\[
y^2=(x-\alpha)^3.
\]
这在$\overline k[x,y]$中不可约：作为$y$的二次多项式，若可约就要求$(x-\alpha)^3$在$\overline k(x)$中为平方，但它在$x=\alpha$处的赋值为3。齐次方程也不含直线$Z=0$为分量，因此$C$几何整。

在$y\ne0$处，$\partial(y^2-x^3-t)/\partial y=2y\ne0$，所以$C$在$k$上光滑。在无穷远只有一点$[0:1:0]$；取$Y=1$的局部方程$z-x^3-tz^3=0$，其关于$z$的偏导数为1，故此点也光滑。剩下的点仅为$q$。局部环$A_{(y)}$是一维诺特整局部环，其极大理想由一个非零元素$y$生成，所以它是正则局部环。故$C$处处正则。

然而$C$在$q$处不光滑。其代数闭包上的点$(\alpha,0)$具有局部方程$y^2=u^3$；极大理想$(u,y)$的$\overline k$-余切空间维数为2，而局部环维数为1，因而该几何局部环不正则。令
\[
W=C\setminus\{q\}.
\]
则$W$是光滑几何整曲线，并含光滑$k$-有理点$[0:1:0]$。

若某个非空开集$U\subset W$是初等的，则按定义它嵌入一条光滑射影几何连通曲线$C'$，且补集有限。$U$在$C'$和$C$中均稠密，故两曲线具有同一函数域。正则射影整曲线由函数域唯一确定：双有理映射在每个离散赋值环处由目标的固有性唯一延拓，所得两方向态射的复合在稠密开集为恒等，由分离性知处处为恒等。因此$C'\simeq C$，与$C$在$q$处不光滑矛盾。也可核对\href{https://stacks.math.columbia.edu/tag/0BY1}{Stacks Project，Tag 0BY1，正则射影曲线与函数域的对应}；该对应说正则，不把不完美域上的正则擅自换成光滑。

若$U\to\Spec(k)$为多重初等分解，每个非空初等态射光滑、满射、相对维数为1；沿分解相加，$\dim U$等于分解中初等态射的数目。由于$\dim U=1$，这个数目只能为1。因此$U$多重初等便已经初等，上述矛盾仍成立。于是$W$的任何点都没有多重初等开邻域。这表明13.2的完美性在其一般性下确实不能删除；结论强于仅指出原证明某一步用了完美性。

\subsection*{四、初段补证三：横截中心及非空边界}

本节先按原证明实际处理的有理点情形，设$x\in X(k)$，并已取经过$x$的$L_0$，使引理13.3的a)、b)、c)成立。记
\[
D=X\cap L_0,\qquad F=(Y\cap L_0)\cup\{x\},\qquad L_0\simeq\mathbf P^N_k.
\]
这里$D$是非空光滑射影曲线，$F$是有限闭集。由$X$纯维$n$以及$L_0$余维$n-1$，交的每个分量维数至少为1；引理b)给维数至多1，所以$D$纯一维。

参数化$L_0$中超平面的对偶空间记为$B=(\mathbf P^N_k)^\vee$。在代数闭包上，给定$z\in D$，使$H\cap D$在$z$处不横截的超平面$H$正是包含射影切直线$\mathbf T_zD$的超平面，其参数维数为$N-2$（$N=1$时为空）。随$z$变化得到的关联闭集在$D\times B$中维数至多$N-1$；由于$D$射影，其投影像闭，且不是维数$N$的整个$B$。其补集即为与$D$横截的超平面参数非空开集。

再排除所有与有限集$F$相交的超平面。对$F_{\overline k}$的每一个几何点，这只是一个真线性闭集；有限并仍不覆盖$B_{\overline k}$。横截与避开的条件都定义在$k$上；因此它们的交是$B$上的非空开集。无限域$k$上射影空间的$k$-点Zariski稠密，故可取$k$-超平面$C\subset L_0$同时满足这两个条件。

由横截性，$D\cap C$是零维光滑$k$-概形；它闭于射影曲线，故有限平展。它还非空：若$C$与某个正维射影分量不相交，该分量就闭嵌入仿射开集$L_0\setminus C$，同时仿射且固有，必为零维，矛盾。设
\[
T=X\cap C=D\cap C.
\]
于是$T$有限平展、非空、与$Y$不交且不含$x$。这正是原稿紧接中心选取后所需要的全部结论。单独“避开$F$”不能代替这里的横截条件。

一般闭点情形不应试图在固定嵌入中要求一个$k$-线性截面经过全部共轭点：这可能根本不存在。13.4(c)的有限几何点方案改为先取共同的$k$-中心$C$，再对每个几何点$x_i$形成$\langle C,x_i\rangle$；其非空开条件、有限点下降及归纳另有本版补证。本节的有理点横截论证只是其中一个局部构造，并未自称替代全部一般点证明。

\subsection*{五、初段补证四：关联族、光滑性、连通性及中心外同构}

\noindent\textbf{1. 参数空间与三个族。}
设$C\subset\mathbf P^r_k$余维$n$，记相应向量子空间为$\widetilde C\subset k^{r+1}$，维数$r-n+1$。包含$C$、余维$n-1$的线性子簇，等价于在$n$维商空间$k^{r+1}/\widetilde C$中选择一维子空间。故
\[
Q(C)\simeq\mathbf P^{n-1}_k.
\]
关联族$X_Q^{(n-1)}$是$X\times_k Q$的闭子概形，因而其投影$p$射影。同样，$Y_Q^{(n-1)}$及$T_Q^{(n-1)}$闭于它。因$Y\cap T=\varnothing$，后二者互不相交。对任意检验$Q$-概形，$T\subset C$确保每一个参数线性子簇都包含$T$的任意拉回，故有概形层面的等式
\[
T_Q^{(n-1)}=T\times_k Q,
\]
不只是几何点集相等。因此$r$有限平展。

\noindent\textbf{2. 相对光滑性的核验。}
在$L_0$的Grassmann参数邻域内，泛线性子簇由$n-1$个独立线性方程局部给出。对$z\in X\cap L_0$，a)保证$z\notin Z$，故$X$在$z$处光滑，维数$n$；b)保证截面在$z$处光滑且维数1。几何扩域后考察余切空间，这意味着$n-1$个方程的微分线性独立。相对Jacobian判据遂表明$p$在$(z,\xi_0)$处光滑且相对维数1。对$z\in Y\cap L_0$同理：$Y$在该处光滑、维数$n-1$，截面维数0，因此$q$在对应点平展。

这一步用到了关联族在光滑环境$X\times Q$或$Y\times Q$中由方程切出的具体结构。一般态射只有一个光滑纤维并不足以推出态射光滑；不能省去这一微分秩核验。光滑位点、平展位点开放；由于$p$、$q$固有，可将这些良好位点的补集的闭像从$Q$中去掉，得到包含$\xi_0$的开集$V$，使$p|_V$光滑、$q|_V$平展。光滑后，相对维数分层为开闭，去掉不等于1的分层的固有像，即可使$p|_V$处处相对维数1。

\noindent\textbf{3. 适用的连通性定理及其完整假设。}
本次已核读Diletta Martinelli、Juan Carlos Naranjo、Gian Pietro Pirola的论文\emph{Connectedness Bertini Theorem via numerical equivalence}，\emph{Advances in Geometry} 17(1) (2017), 31--38，DOI \href{https://doi.org/10.1515/advgeom-2016-0028}{10.1515/advgeom-2016-0028}，特别是定理1.1。作者公开版本见\href{https://diposit.ub.edu/server/api/core/bitstreams/d9e33fe8-dd1c-42a4-a579-a349a203ca25/content}{大学仓储PDF}，另见\href{https://arxiv.org/abs/1412.1978}{arXiv:1412.1978}。该定理适用于任意特征的代数闭域：从不可约射影簇到射影空间的态射，其像与所取线性子空间的维数之和严格大于环境维数时，逆像连通。

对$p$的任一几何参数$\bar\xi$，在其代数闭剩余域上应用该定理于闭浸入$X\hookrightarrow\mathbf P^r$和参数对应的线性子空间$L_{\bar\xi}$。由于$X$几何不可约且射影，假设保持成立；又有
\[
\dim X+\dim L_{\bar\xi}
=n+(r-n+1)=r+1>r.
\]
故$X\cap L_{\bar\xi}$连通。于是$p$的全部几何纤维连通，当然包括缩小到$V$之后的纤维。这是本轮对原稿“Bertini定理”引用的明确核验；并未在本报告重新证明所引用的一般连通性定理。

原稿本句只重提“几何连通”，而该条件单独不足以替代几何不可约。例如取$\mathbf P^4$中的两个射影平面$P_1,P_2$，只在一点$a$相交，并令$X=P_1\cup P_2$。$X$连通、射影、纯二维；一个不经过$a$的一般超平面与两个平面分别相交成不相交的两条直线，故交不连通。这个例子说明必须把原定理中沿用的几何不可约条件保留在应用范围内。

\noindent\textbf{4. 定义13.1的所有条件。}
令$X'=X_Q^{(n-1)}|_V$，
\[
Z'=Y_Q^{(n-1)}|_V\;\amalg\;T_Q^{(n-1)}|_V.
\]
两分量闭且互不相交，所以不交和确为$X'$的闭子概形。$Z'\to V$既固有又平展，故有限平展；其秩局部恒定。$V$非空且为不可约空间$Q\simeq\mathbf P^{n-1}_k$的开集，因而连通，故秩恒定。$T\ne\varnothing$为有限平展$k$-概形，$T\times_k V\to V$为满射；所以$Z'\to V$也满射。此前已经证明$X'\to V$射影、光滑、相对维数1且几何纤维连通，因而$U=X'\setminus Z'\to V$满足初等态射定义的每一条。

\noindent\textbf{5. 中心外同构与正确的邻域包含。}
选择坐标使$C$由$x_0=\cdots=x_{n-1}=0$定义。中心外的线性投影为
\[
\pi_C:\mathbf P^r\setminus C\longrightarrow\mathbf P^{n-1},
\qquad [x_0:\cdots:x_r]\longmapsto[x_0:\cdots:x_{n-1}].
\]
其参数值正是唯一包含$C$与所给点的余维$n-1$线性子簇。对任意检验概形，在某个$x_i$可逆的标准图上，关联条件唯一解出参数坐标为$x_j/x_i$；因此中心外关联族就是$\pi_C$的图。限制到$X$，得到
\[
X_Q^{(n-1)}\setminus T_Q^{(n-1)}\simeq X\setminus T.
\]
进一步删去$Y$便有
\[
X''=X_Q^{(n-1)}\setminus\bigl(Y_Q^{(n-1)}\cup T_Q^{(n-1)}\bigr)
\simeq X\setminus(Y\cup T).
\]
因为$\xi_0\in V$，$U$包含的正是该参数纤维删去边界后的部分
\[
(L_0\cap X)\setminus(Y\cup T).
\]
$x\in L_0$且$x\notin Y\cup T$，故$x\in U$。全纤维$L_0\cap X$包含非空$T$，绝不可能整体包含在$U$内。这里同时核实了法稿$Z_Q$误字的改正以及末句包含关系的改正。

\subsection*{六、一般点衔接与本模块核验限度}

13.2的原陈述允许任意光滑点，下一原段的带点Bertini证明却先使用$x\in P(k)$。单点有理版本不能直接在不完美的残余域上任意改换底域后再下降，也不能把闭点写成$k$-点。主报告对13.4(c)的独立补证采用有限几何点与共同中心：取有限闭点集在代数闭包上的全部点，分别得到好截面开条件；与中心横截$X$且避开$Y$的非空开条件取有限交，再在$k$上选择中心。有限点同时处理的归纳确保后续底空间的所有像点都保留。这个补证正是使原段一般点声明成立的衔接，不能把它归给此处作者已写出的证明。

对于非闭点$x$，在补证一的仿射有限型光滑邻域$W$中，其闭包含有一个在$W$内闭的点$y$；不额外断言$y$在原来的任意预概形$X$内也闭。任意包含$y$的开集也包含其一般化$x$。因此，将有限闭点版应用于$W$所得的开邻域也解决原来的$x$。这一步并不要求$x$的剩余域为有限扩张。取有限点版时，所有给定点先位于共同仿射开集；它们在新底空间的像又为有限个闭点，可取共同仿射开邻域后归纳。

本模块已核验：局部有限表示入口、零维与曲线归约、完美性不可删的反例、射影正规化的适用条件、边界选择、中心横截及非空性、三个关联族的实际结构、相对光滑性、线性截面的连通性定理应用、有限平展边界的定义条件以及中心外同构。带点Bertini的一阶条件及有限几何点下降的完整证明分别在本版后续两个报告模块中；本模块不重写其证明，也不把原笔记早先8.7(c)、8.8等残文宣布为已经恢复。构建、单文件展开与成品页面核验由整合阶段另作真实执行记录；看过源页不等于已编译或看过全书成品。

\subsection*{七、旧D17-M046至D17-M060的逐字存档}

下列旧注作为版本历史保存。它们描述的是旧版制作时的判断，不自动代表本轮结论；尤其“底本”多指英文编辑稿，来源归属以本报告对法稿的逐项核对为准。由独立脚本从实际继承文件的平衡花括号区间提取原文，未用本轮改写文字替换旧注。各原始内容的散列与条数记录在同目录的结构化核查记录中。

\paragraph*{旧注 D17-M046（逐字存档）}
\begin{quote}
底本使用“elementary”和“polyelementary”，这里分别译作“初等”和“多重初等”。按第49页紧随其后的原作者说明，复合允许空族，这解释了恒等态射的情形；这里仍把该说明留作译注，没有冒充定义原段中已有的句子。脚注63未给出Milne书名或SGA IV的具体位置，故不补定引用。
\end{quote}

\paragraph*{旧注 D17-M047（逐字存档）}
\begin{quote}
“多重初等”沿用定义13.1，指结构态射为初等态射的复合。底本定理本句未逐项重述有限型等环境假设，而下文立即使用仿射、拟射影与射影闭包的归约；本译稿不将这些略写解读为对任意无有限性条件的预概形都已给出证明。
\end{quote}

\paragraph*{旧注 D17-M048（逐字存档）}
\begin{quote}
底本此句写“open elementary neighborhood”，而定理陈述是“open polyelementary neighborhoods”，下一段又按维数归纳并构造$U\to V$，故此处正文统一为“多重初等开邻域”，原用词留于此。初等结构态射的纤维为一维，不能直接把一般高维邻域称作在$k$上初等的。句末的字母难辨标记并未因此消除。
\end{quote}

\paragraph*{旧注 D17-M049（逐字存档）}
\begin{quote}
此处保留原作者关于空复合的编写提醒。归纳中还须在$V$上取像点的适当邻域并拉回；底本没有在本段展开，本译文不把省略的归纳拼接冒充为原文已写出的步骤。
\end{quote}

\paragraph*{旧注 D17-M050（逐字存档）}
\begin{quote}
这一段关于曲线模型的写法隐含先在$x$附近作适当开限制，使所讨论的光滑曲线开集置于其正规射影模型中。并且定义13.1要求边界到基底满射，在域上即边界非空；若边界为空，还需在不含$x$处作适当删点。底本这里没有写出这些步骤，故指出略写而不声称已在正文补齐。
\end{quote}

\paragraph*{旧注 D17-M051（逐字存档）}
\begin{quote}
“不改变一个邻域”按局部含义译出，并非断言正规化在整个$X$上是同构。原稿对正规化有限性的引用尚待定位，原脚注65的疑问语气保留。后续构造现已按底本接续收入，但这不等于该依赖链已经独立审定。
\end{quote}

\paragraph*{旧注 D17-M052（逐字存档）}
\begin{quote}
底本把扭曲次数也写为$n>0$，与一直表示$\dim X$的$n$重名；这里改用$d$表示次数，未改变任何数值条件。“取截面”展开为取其非零点开集，以明确基本邻域系的对象。选择包含$Z$的边界是原稿在本段给出的结论，本译稿不添加一个新的选取证明。
\end{quote}

\paragraph*{旧注 D17-M053（逐字存档）}
\begin{quote}
底本的包含关系印为$X\supset Y\subset Z$，与紧前一段$Z\subset Y$及维数限制不符，故订为$X\supset Y\supset Z$。底本重复的$Z$维数限制仍保留。复合符号在底本写作“$\phi_n^{i}$”，并以$n\geq2$为指标；依“替换射影浸入”及第9节的引用，正文明确记为$\phi_d\circ i$，以$d$区别于维数$n$。“any”与“[some?]”相冲突，正文保留不确定性，不断言任意次数均可。
\end{quote}

\paragraph*{旧注 D17-M054（逐字存档）}
\begin{quote}
a）底本写“$\phi=\varnothing$”，这里直接采用空集符号，避免与前段的映射$\phi_d$混淆。
\end{quote}

\paragraph*{旧注 D17-M055（逐字存档）}
\begin{quote}
本句启动的是利用引理的构造，其具体步骤见下段；引理的证明也在后文接续收入，但原稿的依赖与略写仍须分别核查。
\end{quote}

\paragraph*{旧注 D17-M056（逐字存档）}
\begin{quote}
底本两次写$L$，同段的$L_0\cap Y$及随后参数点所代表的子空间明确指向引理中的$L_0$，故统一为$L_0$。但仅仅避开所列有限集，并不足以从字面上推出$X\cap C$平展；还需适当横截选择$C$。此处保留原作者对所选$C$的性质的断言，并明确指出该选择步骤未在本段展开，不把任意满足避开条件的$C$都说成适用。
\end{quote}

\paragraph*{旧注 D17-M057（逐字存档）}
\begin{quote}
底本先写参数点$\xi$，下页续文改写$\xi_0$，此处统一为$\xi_0$。列举第三个逆像时底本只印出带上下标的$Q$而漏掉对象字母；紧随其后的$T_Q^{n-1}\simeq T\times_k Q$及边界构造确定它是$T$的逆像。括号中的上标$(n-1)$表示余维，排版整理不改变原含义。
\end{quote}

\paragraph*{旧注 D17-M058（逐字存档）}
\begin{quote}
这一连通性结论按原稿引用保留。底本并未在此写明所用Bertini定理的完整假设或引用位置，本译稿不以“$X$几何连通”单独替代所需的线性截面连通性论证。
\end{quote}

\paragraph*{旧注 D17-M059（逐字存档）}
\begin{quote}
底本原句为“an open neighborhood $U$ of $Q$ in $X$”，但紧接着把族限制到$V$，上下文始终是在参数点附近缩小基底，故正文订为“$\xi_0$在$Q$中的开邻域$V$”；原“fiber of $Q$”同理订为$\xi_0$处的纤维。定义13.1所需射影性可由原射影关联族限制得到；边界的满射性和恒定次数还须使用非空$T$及必要的缩小。本段原文没有把这些核对全部写出。
\end{quote}

\paragraph*{旧注 D17-M060（逐字存档）}
\begin{quote}
底本在$X''$的公式中写$Z_Q^{n-1}$，但随即说明它是$X-Y-T$的逆像，且本段移去的第二个边界一直是$T$，故订为$T_Q^{(n-1)}$。底本又写$U$包含整个$L_0\cap X$；因$U$已删去边界，这应为删去$Y\cup T$后的部分，正文据构造修正。最后“因而包含$x$”还使用$x\in L_0$：底本引理13.3的陈述漏明这一点，下一段证明却明确取经过$x$的线性子簇。此项前后依赖在这里指出，不能只凭前述引理字面推出。
\end{quote}

